\section{The Enrichment Layer: From Score to Behavioral Corpus}
\label{sec:enrichment}

The scorecard of Section~\ref{sec:grading} answers \emph{did the seller run a good
deal}. It does not, by itself, answer the questions a researcher or an RL engineer
asks next: \emph{where} did the deal turn, \emph{how} did the buyer's posture move,
\emph{which} sentence created urgency, and \emph{whether} the model held value
because it understood the buyer or merely recited a script. VEB therefore ships a
second, orthogonal layer on top of every graded trajectory---an
\emph{enrichment} pass that converts each transcript into a densely labeled,
machine-readable behavioral record. Where SQS is the benchmark's \emph{verdict},
the enrichment layer is its \emph{instrument}: it is what turns \resEnrichRows{}
graded episodes into a reusable research and post-training corpus rather than a
leaderboard snapshot.

\subsection{Two feature families}
Each episode is enriched with \numDetFeatures{} \textbf{deterministic} features and
\numAiFeatures{} \textbf{behavioral} features (15 AI-judged plus the two
deterministic divergence signals, which ride in the same record with the same
provenance schema). Deterministic features are computed
directly from the transcript and the internal channel with no model call---turn
counts, cadence and latency between touches, discount trajectory, EB-meeting
timing, quote-fidelity flags, and a pair of public/private \emph{divergence}
signals that fire when the seller's private plan contradicts its public
posture. These are exact and free, and they anchor the layer: a reader who
distrusts any AI feature can fall back on them.

The AI-judged features capture what deterministic parsing cannot---the read of a
room. They are organized into \numEnrichPasses{} passes, each a separate prompt
over the same transcript so that a single pass failure never corrupts the others:

\begin{itemize}[leftmargin=1.4em,itemsep=2pt,topsep=2pt]
  \item \textbf{People} (5 features): buyer trust read, seller confidence, buyer
  frustration peak (lower is better), buyer posture
  (\texttt{guarded}/\texttt{engaged}/\texttt{evaluating}/\texttt{committing}/\texttt{dark}),
  and the seller's theory-of-mind gap (how far its model of the buyer diverges
  from the buyer's actual private state; lower is better).
  \item \textbf{Message} (5 features): value specificity, question quality,
  linguistic clarity, discovery-vs-pitch balance
  (\texttt{discovery\_led}/\texttt{balanced}/\texttt{pitch\_led}), and
  objection-handling quality.
  \item \textbf{Play} (5 features): deal-read accuracy, value defended
  (\texttt{held}/\texttt{traded}/\texttt{gifted}), the pivotal-turn index (the
  single turn most responsible for the outcome), counterfactual lift (how much
  that pivotal move changed the expected result), and urgency created.
\end{itemize}

\subsection{Cite-or-zero}
\label{sec:cite-or-zero}
The enrichment layer inherits the benchmark's central discipline: \emph{every
AI-judged feature must cite the transcript turns that justify it, or its value is
voided to zero}. A feature is not permitted to express a confident number with no
evidence behind it; an uncited claim is treated as no claim. This makes the entire
labeled corpus auditable in the same way the scorecard is---any downstream
consumer can click a feature and read the exact turns the label rests on---and it
structurally suppresses the hallucinated-insight failure mode that makes naive
``LLM-annotates-transcript'' pipelines untrustworthy for training data.

\subsection{Ensemble scoring}
Each AI feature is scored by a \numEnrichSeats{}-seat cross-family ensemble rather
than a single model, and the released value is the seat consensus with a retained
per-seat breakdown and a calibrated confidence. Using an ensemble here serves the
same anti-self-preference purpose as the SQS panel (Section~\ref{sec:panel}) but
at the finer granularity of individual behavioral judgments, where a lone model's
idiosyncratic reading of ``urgency'' or ``posture'' would otherwise pass through
unchecked. The production ensemble is a deliberately \emph{cost-reduced} panel; we
justify that choice empirically in Section~\ref{sec:panel-calibration}, where we
show it tracks a frontier reference panel at $r=\resPanelPearson{}$ on real data
before committing it to the full sweep.

\subsection{What the layer is for}
The enrichment layer is the bridge between VEB-as-evaluation and VEB-as-training
environment. For analysis, it lets us decompose \emph{why} a model ranks where it
does---separating models that convert a methodology into behavior from those that
convert it only into vocabulary (Section~\ref{sec:results}). For post-training, the
seed-matched trajectories carry per-turn, evidence-cited behavioral labels, so a
reward model can be shaped on \emph{how} a deal was won---trust earned, value held,
the buyer's own words elicited---rather than on the terminal win bit alone. We
release the full feature catalog, the deterministic-vs-AI provenance of every
column, and the per-seat ensemble outputs with the dataset.
